% Primera lectura: puntos y coordenadas, transformar, rotar, reflejar, escalar, afinidades, componer con centros distintos
\section{El punto y sus coordenadas no son lo mismo}\label{sec:punto-coordenadas}

Un punto geométrico existe independientemente de la manera elegida para describirlo. Sus
coordenadas son una lista de números obtenida después de elegir un origen, unos ejes, una
orientación y una unidad. Si cambia el sistema de coordenadas, cambian los números, pero no
necesariamente cambia el punto.

\begin{idea}[Idea conceptual]
La materia estudia simultáneamente dos problemas distintos:
\begin{itemize}
  \item mover o deformar objetos geométricos;
  \item cambiar la manera de asignarles coordenadas.
\end{itemize}
Las fórmulas pueden parecer similares, pero la interpretación activa o pasiva determina qué
está cambiando realmente.
\end{idea}

\paragraph{Primera traducción algebraica.}
Una posición puede representarse mediante un vector \(x\). Cambiar de descripción significa
cambiar sus coordenadas; transformar activamente el objeto significa producir un nuevo punto
\(x'=F(x)\). La fórmula es breve, pero su interpretación depende de qué permanece fijo:
el objeto, los ejes o ambos.

\begin{idea}[En palabras del profesor]
``La rotación genera para cada punto una trayectoria o línea de los lugares posibles a
los que caer. [\ldots] Estas trayectorias posibles son lo que conocemos normalmente como
líneas coordenadas'' (resumen de clase 1, cursada 2025).
\end{idea}

\section{Qué significa transformar una figura}\label{sec:transformar}

Una transformación asigna a cada punto una imagen. Para comprenderla geométricamente conviene
imaginar una figura sencilla, como un triángulo, y observar:
\begin{itemize}
  \item hacia dónde se desplaza cada vértice;
  \item si los lados conservan longitud;
  \item si los ángulos se mantienen;
  \item si las rectas paralelas continúan siendo paralelas;
  \item si existe un punto, una recta o una dirección que permanezca fija.
\end{itemize}

\begin{center}
\begin{tikzpicture}[scale=.88]
  \coordinate (a) at (1,1);
  \coordinate (b) at (3.4,1.2);
  \coordinate (c) at (2,3.2);
  \coordinate (ap) at (7,1.4);
  \coordinate (bp) at (10.2,1.9);
  \coordinate (cp) at (8.2,4);
  \draw[teal,line width=1.1pt] (a)--(b)--(c)--cycle;
  \draw[rust,line width=1.1pt] (ap)--(bp)--(cp)--cycle;
  \foreach \p/\q in {a/ap,b/bp,c/cp}
    \draw[auxline,-{Latex[length=2mm]}] (\p)--(\q);
  \node[labelback,text=teal] at (2.1,.55) {figura original};
  \node[labelback,text=rust] at (8.7,1.0) {figura imagen};
\end{tikzpicture}
\end{center}

La matriz y el vector de traslación son una descripción compacta de ese movimiento. Antes de
operar con ellos, debe poder anticiparse cualitativamente la figura resultante.

\paragraph{Primera traducción algebraica.}
\[
  F(x)=Ax+t.
\]
Aquí \(A\) transforma direcciones y formas, mientras que \(t\) desplaza toda la figura sin
deformarla. Las columnas de \(A\) muestran dónde terminan los vectores unitarios de los ejes.

\begin{idea}[En palabras del profesor]
``Hay transformaciones que se definen completamente por qué propiedades de los objetos
transformados mantienen'' (resumen de clase 1, rúbrica, cursada 2025).
\end{idea}

\section{Rotar: conservar distancias alrededor de un centro}\label{sec:rotar}

Una rotación mueve cada punto sobre una circunferencia cuyo centro es el centro de rotación.
Por eso conserva:
\begin{itemize}
  \item la distancia de cada punto al centro;
  \item las distancias entre puntos;
  \item los ángulos y la forma;
  \item la orientación.
\end{itemize}

En dos dimensiones basta un centro y un ángulo. En tres dimensiones también se necesita un
eje: la dirección del eje permanece fija y todo plano perpendicular a él gira como un plano
bidimensional.

\paragraph{Primera traducción algebraica.}
\[
  x'=c+R(\theta)(x-c).
\]
El vector \(x-c\) expresa la posición vista desde el centro; \(R(\theta)\) lo gira, y sumar
\(c\) devuelve el resultado al sistema original.

\begin{resultado}[Lectura inmediata]
Si una transformación 3D es una rotación, debe existir una dirección que no cambie. Esa
dirección es el eje y más adelante aparecerá algebraicamente como un autovector de autovalor
\(1\).
\end{resultado}

\begin{idea}[En palabras del profesor]
``Una rotación implica dejar quieto un punto, mover otro, y en base a esos acomodar todos
los demás. [\ldots] Se podían acomodar los demás porque había una estructura de algo que
se conservaba (distancias y ángulos)'' (resumen de clase 1, cursada 2025).
\end{idea}

\section{Reflejar y componer}\label{sec:reflejar}

Una reflexión invierte una figura respecto de un centro o de una recta. Aplicada dos veces
respecto del mismo elemento devuelve la figura original. Esta propiedad se llama involución.

La riqueza geométrica aparece al combinar reflexiones distintas:
\begin{itemize}
  \item dos reflexiones puntuales producen una traslación;
  \item dos reflexiones en rectas paralelas producen una traslación perpendicular a ellas;
  \item dos reflexiones en rectas secantes producen una rotación alrededor de su intersección;
  \item tres reflexiones puntuales vuelven a producir una reflexión puntual y construyen un
  paralelogramo.
\end{itemize}

\paragraph{Primera traducción algebraica.}
\[
  \rho_b(x)=2b-x.
\]
La diferencia \(\rho_b(x)-b=-(x-b)\) indica que la imagen conserva la distancia a \(b\),
pero queda en la dirección opuesta. Al componer dos reflexiones, los términos que dependen de
\(x\) recuperan su signo y queda un desplazamiento constante.

\begin{idea}[Qué debe observarse]
Cada reflexión individual cambia la orientación. Dos reflexiones vuelven a preservarla. La
posición relativa de los centros o de los ejes decide si el resultado se interpreta como
traslación o como rotación.
\end{idea}

\begin{idea}[En palabras del profesor]
``Vimos que con reflexiones podíamos construir rotaciones (con un argumento coloquial
sobre que un par de reflexiones sucesivas no altera las distancias, pero tiene un punto
fijo, así que tiene que ser una rotación)'' (resumen de clase 3, cursada 2025).
\end{idea}

\section{Escalar: homotecias y similaridades}\label{sec:escalar}

Una homotecia cambia el tamaño respecto de un centro. Cada punto permanece sobre la recta que
lo une con ese centro. Si la razón es positiva conserva el semirayo; si es negativa pasa al
semirayo opuesto.

Una similaridad combina escala uniforme, giro y, eventualmente, reflexión o traslación. La
figura puede cambiar de tamaño y posición, pero conserva sus ángulos y sus proporciones. Dos
triángulos relacionados por una similaridad tienen la misma forma.

\paragraph{Primera traducción algebraica.}
\[
  x'=o+k(x-o).
\]
El término \(x-o\) mira el punto desde el centro. Multiplicarlo por \(k\) cambia su distancia
y, si \(k<0\), también invierte el semirayo.

\begin{idea}[En palabras del profesor]
``Así como la rotación se expresa en complejos también se expresa la escala, que en
matrices se representa con una matriz múltiplo de la identidad'' (resumen de clase 1, cursada 2025).
\end{idea}

\section{Afinidades: conservar rectas y paralelismo}\label{sec:afinidades-primera}

Una afinidad permite deformaciones más generales. Un cuadrado puede convertirse en
paralelogramo y un círculo en elipse, pero las rectas siguen siendo rectas y las paralelas
siguen siendo paralelas cuando la transformación es invertible.

\begin{center}
\begin{tikzpicture}[scale=.86]
  \draw[teal,line width=1.1pt] (1,1) rectangle (3.7,3.7);
  \draw[rust,line width=1.1pt] (7,1)--(10.2,1.5)--(9.1,4.2)--(5.9,3.7)--cycle;
  \draw[-{Latex[length=2.8mm]},tealdark,line width=1pt]
    (4.3,2.4)--(5.4,2.4);
  \node[labelback,text=teal] at (2.35,.55) {cuadrado};
  \node[labelback,text=rust] at (8.05,.55) {paralelogramo};
\end{tikzpicture}
\end{center}

Los puntos fijos muestran dónde se ``ancla'' la transformación. Las direcciones propias
muestran líneas que conservan su dirección aunque puedan estirarse, contraerse o invertirse.
Las coordenadas homogéneas permiten tratar la traslación dentro del mismo lenguaje matricial.

\paragraph{Primera traducción algebraica.}
\[
  F(x)=Ax+t.
\]
La parte \(A\) deforma el paralelogramo de la base y \(t\) transporta su origen. Si \(A\) es
invertible, ninguna dirección completa colapsa en un punto.

\begin{idea}[En palabras del profesor]
``Una afinidad conserva paralelas las paralelas'' o ``una afinidad conserva rectas las
rectas y proporcionales las distancias sobre una recta'' (resumen de clase 3, cursada 2025).
\end{idea}

\section{Componer transformaciones con centros distintos}\label{sec:componer-centros}

La clase práctica combina dos operaciones:
\begin{enumerate}
  \item una rotación de \(45^\circ\) alrededor de un punto \(c\);
  \item una escala uniforme de razón \(k=0.5\) alrededor de otro punto \(d\).
\end{enumerate}
La primera cambia la orientación sin cambiar el tamaño. La segunda reduce todas las distancias
al centro \(d\) a la mitad. Como los centros son distintos, la transformación final no puede
entenderse mirando solamente una de las dos operaciones.

\paragraph{Primera traducción algebraica.}
\[
F_1(x)=c+R_{45^\circ}(x-c),
\qquad
F_2(y)=d+0.5(y-d).
\]
La composición es
\[
F(x)=F_2\bigl(F_1(x)\bigr).
\]
La expresión situada más a la derecha actúa primero: \(F_1\) rota y después \(F_2\) escala.
Al desarrollar,
\[
F(x)=0.5R_{45^\circ}x+t,
\]
donde el vector \(t\) reúne el efecto de trasladar hacia ambos centros y regresar.

\begin{idea}[Qué clase de transformación resulta]
La parte lineal \(0.5R_{45^\circ}\) gira \(45^\circ\) y reduce a la mitad. Por tanto, la
composición completa es una similaridad directa contractiva. Posee un único punto fijo, que
actúa como centro global de la transformación compuesta aunque no coincida, en general, con
ninguno de los dos centros originales.
\end{idea}

\paragraph{Movimiento discreto y trayectoria imaginada.}
Matemáticamente, \(F\) lleva cada punto directamente a su imagen. La trayectoria ``primero
un arco y después un segmento recto'' es una interpolación útil para comprender las dos
etapas, pero no forma parte de la definición de la transformación compuesta. Otra
interpolación podría unir los mismos extremos por un camino diferente.

\begin{idea}[En palabras del profesor]
``Experimentamos qué pasaba cuando aplicábamos rotaciones y traslaciones sucesivamente a
un punto (formar círculos) y tratamos de pensar por qué esto era inevitable, con un
argumento nuevamente sobre el punto fijo de la transformación'' (resumen de clase 3, cursada 2025).
\end{idea}

